\section{Linear path stability for planar graphs and bounded genus}
\label{sec:path-phase-two}

This section gives quantitative bounds for planar graphs, using the small-cut lemma and a reduction adapted to one path.
The same argument gives a bound with an additive genus term.
A separator argument also proves linear stability for every fixed proper minor-closed class.
These results are not formalized in Lean.

Write $P_{uv}(N)$ for the number of simple paths between distinct vertices $u,v$.
Write $b=\beta(N)$.
The orientable genus of a planar graph is zero.

\begin{theorem}[Linear stability with a genus term]
\label{thm:path-genus}
Let $N$ be a connected simple graph that embeds in an orientable surface of genus at most $g$.
Let $A$ be a $u$--$v$ path of length $\ell$.
Suppose that
\[
 P_{uv}(N)\ge(1/2+\varepsilon)2^b,
 \qquad 0<\varepsilon\le1/2.
\]
For every integer $L\ge6$,
\begin{equation}\label{eq:path-genus-explicit}
 (L-5)b
 \le 2(L+3)\ell+2(L+1)g+(L+1)
 +\frac{3L(L+1)(L-1)2^{L-2}}{\varepsilon}.
\end{equation}
In particular,
\[
 b\le10\ell+8g+4+\frac{16128}{\varepsilon}.
\]
For every fixed $\eta>0$, a constant $C_\eta$ gives the stronger asymptotic form
\[
 b\le(2+\eta)(\ell+g)+\frac{C_\eta}{\varepsilon}.
\]
\end{theorem}

\begin{corollary}[The planar case]
\label{cor:path-planar}
Under the same density hypothesis, if $N$ is planar, then
\[
 \ell\ge\frac{b-4-16128/\varepsilon}{10}.
\]
For all ranks, the positive linear bound
\[
 \ell\ge\frac{\varepsilon}{14\varepsilon+16128}\,b
\]
also holds.
For each fixed $\varepsilon>0$ and each $\eta>0$, one has
$\ell\ge b/(2+\eta)-O_{\eta,\varepsilon}(1)$.
\end{corollary}

\begin{proof}[Derivation of the corollary]
Set $g=0$ and $L=7$ in Theorem~\ref{thm:path-genus}.
Since $\ell\ge1$, its upper bound for $b$ is at most
$(14+16128/\varepsilon)\ell$.
The last assertion uses the theorem's final inequality.
\end{proof}

\subsection{A reduction that preserves a short marked path}

\begin{lemma}\label{lem:path-planar-split}
Let $N,A,u,v$ be as above, without a path-density assumption.
There is a connected multigraph $\Gamma$ on the same surface and a terminal path $A'$ such that
\begin{enumerate}
\item $\beta(\Gamma)=b$ and $P_{uv}(N)\le P_{u'v'}(\Gamma)$;
\item every vertex outside $A'$ has degree three;
\item every internal vertex of $A'$ has degree two or three, and its endpoints have degree one or two;
\item $a:=|V(A')|\le2\ell$ and $|V(\Gamma)|\le2b+a$.
\end{enumerate}
\end{lemma}

\begin{proof}
First delete leaves other than the terminals, until no such leaves remain.
This preserves the rank and the terminal paths.
It preserves $A$ as well.

Replace each vertex by a tree inside a small disk in the fixed surface embedding.
The external edges retain their cyclic order around that disk.
At a vertex outside $A$, use a tree whose vertices have degree three, counting external edges.
A degree-two vertex can instead remain unsplit at this stage.

Consider an internal vertex of $A$.
Its two path edges divide the other incident edges into two intervals in the cyclic order.
If both intervals are nonempty, use two marked vertices joined by one edge.
Attach one path edge to each marked vertex.
Attach the two intervals through two rooted trees, one at each marked vertex.
Each marked vertex has degree three.
If only one interval is nonempty, one marked vertex suffices, with both path edges and one rooted tree.
If both intervals are empty, retain the degree-two vertex.
These trees lie in the disk and respect its boundary order.
Their unmarked vertices have degree three.

At each terminal, use one marked vertex.
Its incident path edge is one branch.
Attach all its remaining incident edges through one rooted tree in the other branch.
This marked vertex has degree two, or degree one when there are no remaining edges.
All unmarked vertices of the rooted tree have degree three.
For a rooted tree with only one external edge, use that edge directly.

Replacing a vertex by a tree does not change the rank.
Each original terminal path has a unique lift through these trees.
That lift is simple, since the original path visits each original vertex at most once.
Distinct original paths have distinct lifts.
There can be additional paths after splitting, so the required comparison is an inequality.

The lifted path uses at most two vertices for each original internal path vertex and one for each endpoint.
Thus $a\le2(\ell-1)+2=2\ell$.
Suppress every degree-two vertex outside the lifted path.
This preserves the rank, the embedding, and the path count.
It gives the required multigraph $\Gamma$ and path $A'$.

Write $n=|V(\Gamma)|$.
Every vertex outside $A'$ has degree three.
The internal marked vertices have degree at least two, and the two endpoints have degree at least one.
Therefore
\[
 2(n+b-1)=2|E(\Gamma)|\ge3(n-a)+2(a-2)+2=3n-a-2.
\]
This gives $n\le2b+a$.
\end{proof}

The use of two marked vertices is important.
For arbitrary cyclic orders, grouping both path edges at one marked cubic vertex need not preserve the embedding.

\subsection{Short faces supply disjoint cyclic regions}

\begin{lemma}\label{lem:path-face-packing}
Let $\Gamma,A',a,b,g$ satisfy Lemma~\ref{lem:path-planar-split}.
For each integer $L\ge6$, there are $k$ pairwise vertex-disjoint cycles outside $V(A')$, each of length at most $L$, with
\begin{equation}\label{eq:path-packing}
 k\ge
 \frac{(L-5)b-(L+3)a-2(L+1)g-(L+1)}{3L(L+1)}.
\end{equation}
A negative right side imposes no requirement.
Each chosen vertex set is a cyclic region whose edge boundary in $\Gamma$ has size at most $L$.
\end{lemma}

\begin{proof}
Let $J=\Gamma-V(A')$.
The marked vertices have degree at most three and contain the $a-1$ edges of $A'$.
Thus at most $3a-(a-1)=2a+1$ edges disappear.
If $J$ is nonempty, this gives
\[
 \beta(J)\ge b-a-1.
\]
The inequality also holds when $J$ is empty.
In that case subcubicity gives $b\le a/2+1\le a+1$.
The total number of face-boundary edge occurrences is at most
\[
 2|E(J)|\le2|E(\Gamma)|
 =2(|V(\Gamma)|+b-1)\le6b+2a-2.
\]

In the planar case there are exactly $\beta(J)$ bounded faces.
For a surface of genus $g$, take a regular neighborhood of each component and cap its boundary curves.
This gives cellular embeddings with total genus at most $g$.
Discard one face from each component.
Euler's formula leaves at least $\beta(J)-2g$ faces.
Each remaining face has a boundary containing a cycle.
Tree components contribute no remaining faces.
These statements permit repeated boundary edges and parallel edges.

Call a retained face short when its total boundary length is at most $L$.
Every other retained face has at least $L+1$ boundary occurrences.
The number $f_L$ of short faces therefore satisfies
\[
 \begin{split}
 f_L&\ge b-a-1-2g-\frac{6b+2a}{L+1}\\
 &=\frac{(L-5)b-(L+3)a-2(L+1)g-(L+1)}{L+1}.
 \end{split}
\]
If this bound is positive, choose a cycle in a short face boundary.
Delete from the list every short face incident to a vertex of that cycle.
Repeat until the list is empty.

Every chosen cycle has at most $L$ vertices.
Each such vertex is incident to at most three faces, since its degree in $J$ is at most three.
Thus each choice removes at most $3L$ faces from the list.
The selected cycles are vertex-disjoint, so $k\ge f_L/(3L)$.
This proves \eqref{eq:path-packing}.

Every selected vertex has degree three in $\Gamma$ and uses two incidences in its cycle.
Its cyclic region therefore has at most $|C|\le L$ boundary edges.
\end{proof}

\subsection{The probability estimate}

We use a slightly sharper constant in the manuscript's small-cut lemma.
For disjoint proper cyclic regions with boundary at most $L$, the cut formula gives $p_i\ge2^{-L}$.
The exceptional-partner proof bounds the partners with at least three links by $L-2$.
Consequently, for $Z=\sum_i\mathbf1_{E_i}$ and $\mu=\mathbb EZ$,
\[
 \mu\ge k2^{-L},\qquad
 \mathbb EZ^2\le2\mu^2+(L-1)\mu.
\]
The same Cauchy--Schwarz argument gives
\begin{equation}\label{eq:path-small-cuts}
 \Pr\bigl(\text{the even subgraph is a forest in every region}\bigr)
 \le\frac12+\frac{(L-1)2^{L-2}}{k}.
\end{equation}
Here $k\ge1$, and every region in our application is proper because it avoids $A'$.

\begin{proof}[Proof of Theorem~\ref{thm:path-genus}]
Apply Lemma~\ref{lem:path-planar-split}.
Add a new two-edge terminal path to $\Gamma$, and let $H$ be the resulting graph.
Let $W$ be its union with $A'$.
The path injection from the first research note gives
\[
 \frac12+\varepsilon
 \le\frac{P_{uv}(N)}{2^b}
 \le\frac{P_{u'v'}(\Gamma)}{2^b}
 \le p(H,W).
\]
The genus of $H$ is irrelevant.
The regions supplied by Lemma~\ref{lem:path-face-packing} lie in the embedded graph $\Gamma-V(A')$.
Their cuts are unchanged when the new path is added.
On the forest event outside $W$, the even subgraph is a forest inside each such region.
Thus \eqref{eq:path-small-cuts} yields
\[
 k\le\frac{(L-1)2^{L-2}}{\varepsilon}.
\]
Combine this with \eqref{eq:path-packing} and $a\le2\ell$.
If the right side of \eqref{eq:path-packing} is nonpositive, the desired inequality already holds.
Otherwise the displayed upper bound on $k$ proves \eqref{eq:path-genus-explicit}.

For $L=7$, division by $L-5=2$ gives the stated numerical bound.
For the last assertion, choose $L$ so large that $16/(L-5)\le\eta$.
Then
\[
 \frac{2(L+3)}{L-5}\le2+\eta,
 \qquad
 \frac{2(L+1)}{L-5}\le2+\eta.
\]
Absorb the remaining constant into $C_\eta/\varepsilon$.
\end{proof}

\subsection{Restrictions on a possible general proof}

If the path family has more than half of its affine join space, every even edge set is the difference of two paths.
Indeed, a set and each of its translates must intersect when both have size greater than half the ambient space.
This is necessary, but it does not supply a quantitative length bound.
The chain-with-a-shortcut examples satisfy this condition at arbitrarily large rank.

There is also a useful local restriction.
In a two-connected block, the restriction of a uniform even subgraph to the edges at a vertex is uniform over all even patterns.
Any such pattern extends by a join in the connected graph obtained after deleting that vertex.
For a terminal join, the terminal patterns are instead uniform over all odd patterns.
Thus path density above one half forces internal block degrees at most five and terminal block degrees at most three.
The relevant local probabilities are
\[
 \frac{1+\binom d2}{2^{d-1}}\quad\text{and}\quad
 \frac{d}{2^{d-1}},
\]
which are at most one half for $d\ge6$ and $d\ge4$, respectively.
Ranks and normalized path counts factor over the blocks on the terminal route.
Therefore this observation applies separately to each such block.
It still permits large cubic blocks.

The tempting general bound $P_{uv}\le2^{b-1}+2^\ell-1$ is false.
Start with $K_{3,3}$ on parts $\{0,1,2\}$ and $\{3,4,5\}$, and delete the edge $03$.
Add a vertex $6$ and edges $63,64$.
For terminals $0,6$, the resulting simple graph has rank four and distance two.
It has eight paths ending through $36$ and five ending through $46$, for thirteen paths in total.
The proposed upper bound would be $8+4-1=11$.
The following family rules out that broader form as well.

\begin{proposition}[Unbounded excess at distance two]
\label{prop:path-distance-two-obstruction}
For each integer $b\ge3$, there is a connected simple planar graph of rank $b$, terminal distance two, and
\[
 P_{uv}=2^{b-1}+b.
\]
Consequently, there is no bound $P_{uv}\le2^{b-1}+f(\ell)$ for all connected simple graphs, for any finite-valued function $f$.
In particular, $2^{C\ell}$ cannot serve as an additive error for any fixed $C$.
\end{proposition}

\begin{proof}
Take the ladder on columns $(x_i,y_i)$, $0\le i\le b-1$.
Its edges are the two rails $x_ix_{i+1},y_iy_{i+1}$ and all rungs $x_iy_i$.
It has rank $b-1$.
Each $x_0$--$x_{b-1}$ path crosses every two-edge cut between consecutive columns exactly once.
The choices of upper or lower rail at the $b-1$ cuts determine the path uniquely.
All $2^{b-1}$ choices give simple paths.

Add the edge $x_0x_{b-1}$, and call the resulting planar graph $G$.
Its rank is $b$.
The old cycles are rectangles between two columns.
Exactly $b-1$ of these contain $x_0$.
Every new cycle uses the added edge and corresponds to one of the $2^{b-1}$ terminal paths.
Hence
\[
 c_{x_0}(G)=2^{b-1}+b-1.
\]

The vertex $x_0$ has degree three.
Replace it by three vertices $p_0,p_1,p_2$ joined as a path.
Assign one of its three former incident edges to each new vertex.
This local tree replacement preserves planarity and rank.
Use $p_0,p_2$ as the terminals.
Their distance is two.

The path $p_0p_1p_2$ is one terminal path.
Every other terminal path has exactly one segment outside these three vertices.
Its endpoints are two of the three ports.
Contracting the replacement tree turns that segment into a cycle through $x_0$.
Conversely, every such cycle gives one terminal path, using zero or one of the replacement edges as needed.
Thus the terminal path count is $1+c_{x_0}(G)=2^{b-1}+b$.
\end{proof}

The normalized gap in this family is $b/2^b$, which tends to zero.
It therefore does not contradict fixed-gap linear stability.
It shows that a successful general error bound must retain some rank dependence.

The planar proof isolates its extra input: many short faces give many bounded-cut cyclic regions.
General high-girth cubic graphs have no corresponding face count.
The argument above therefore does not settle the unrestricted linear problem.
Theorem~\ref{thm:general-linear-path-stability} settles that problem through localization and shortest-path layers.

\section{Linear path stability with a fixed excluded minor}
\label{sec:path-minor-stability}

The following extension uses the Alon--Seymour--Thomas separator theorem.
It does not require vertex splitting, which can change an excluded-minor condition.

\begin{theorem}[Fixed excluded minor]\label{thm:path-excluded-minor}
Fix an integer $h\ge2$ and a real number $\eta>0$.
There is a constant $C_{h,\eta}$ with the following property.
Let $N$ be a connected simple graph with no $K_h$ minor and rank $b$.
If
\[
 P_{uv}(N)\ge(1/2+\varepsilon)2^b,
 \qquad 0<\varepsilon\le1/2,
\]
then every $u$--$v$ path of length $\ell$ satisfies
\begin{equation}\label{eq:path-minor-linear}
 b\le(3+\eta)\ell+\frac{C_{h,\eta}}{\varepsilon}.
\end{equation}
In particular, for fixed $h$ and $\varepsilon$, every such path has length $\Omega_{h,\varepsilon}(b)$.
The conclusion applies to every fixed proper minor-closed class.
\end{theorem}

\subsection{Degree and rank restrictions from path density}

\begin{lemma}\label{lem:path-route-reduction}
Suppose that $P_{uv}(N)>2^{b-1}$, and fix a terminal path $A$ of length $\ell$.
There is a connected loopless multigraph $\Gamma$ with the following properties.
\begin{enumerate}
\item Its underlying simple graph is a minor of $N$.
\item Its rank and terminal path count equal those of $N$, and it retains $A$.
\item Every vertex outside $A$ has degree between three and five.
\item It has at most $2b+\ell+1$ vertices, and $\beta(\Gamma-V(A))\ge b-3\ell$.
\end{enumerate}
Degrees count parallel edges separately.
\end{lemma}

\begin{proof}
The blocks on the terminal route have ranks $b_i$ and terminal path counts $P_i$.
Let $s$ be the sum of the ranks of all other blocks.
The affine join bound gives $P_i\le2^{b_i}$, and
\[
 \frac{P_{uv}(N)}{2^b}=2^{-s}\prod_i\frac{P_i}{2^{b_i}}.
\]
The strict density assumption implies $s=0$ and $P_i/2^{b_i}>1/2$ for every route block.
Delete all edges outside the terminal route.
These edges belong to trees, so their deletion preserves both the rank and the terminal paths.

Consider a two-connected route block $B$, with its two route terminals.
Deleting any vertex leaves a connected graph.
Thus every incident edge pattern with the required parity extends to a terminal join.
The restriction of a uniform join to the incident edges is consequently uniform among those patterns.
At an internal vertex of degree $d$, a simple terminal path uses zero or two incident edges.
At a terminal, it uses exactly one.
The corresponding probability bounds are
\[
 \frac{1+\binom d2}{2^{d-1}},\qquad \frac{d}{2^{d-1}}.
\]
The first is at most one half for $d\ge6$, and the second for $d\ge4$.
Therefore internal block degrees are at most five and terminal block degrees are at most three.
Bridge blocks also satisfy these bounds.

Every articulation on the terminal route lies on $A$.
Consequently, all vertices outside $A$ have degree at most five.
Suppress degree-two vertices outside $A$ until none remain.
Suppression preserves the rank and gives a bijection on terminal paths.
It is an edge contraction, with parallel edges retained.
It therefore preserves the required minor exclusion for the underlying simple graph.
No loop arises, since each cyclic route block retains its two distinct terminals.
The resulting graph is $\Gamma$.

Write $n=|V(\Gamma)|$ and $a=|V(A)|=\ell+1$.
The degree lower bounds give
\[
 2(n+b-1)\ge3(n-a)+2(a-2)+2=3n-a-2.
\]
Hence $n\le2b+\ell+1$.

For the rank bound, consider one route block after suppression.
Let its part of $A$ have length $t\ge1$.
The sum of degrees at its $t+1$ marked vertices is at most $5(t-1)+6=5t+1$.
These vertices contain the $t$ path edges.
At most $4t+1$ edges therefore disappear when the marked vertices are deleted.
If an unmarked vertex remains, the surviving rank is at least
\[
 \beta(B)+(t+1)-(4t+1)=\beta(B)-3t.
\]
If no unmarked vertex remains, then
$\beta(B)\le(5t+1)/2-t\le3t$, giving the same conclusion.
After the marked vertices are removed, the route blocks become vertex-disjoint.
Their ranks add, and their marked path lengths sum to $\ell$.
Thus $\beta(\Gamma-V(A))\ge b-3\ell$.
\end{proof}

\subsection{Bounded components from separators}

We use the following form of the Alon--Seymour--Thomas theorem.
An $m$-vertex graph with no $K_h$ minor has a vertex set of size at most $h^{3/2}\sqrt m$.
After its deletion, every component has at most $m/2$ vertices.
This is statement (1.4) of their original paper.\footnote{N. Alon, P. Seymour, and R. Thomas,
\emph{A separator theorem for nonplanar graphs}, Journal of the American Mathematical Society 3 (1990), 801--808.
The original paper is available at \url{https://thomas.math.gatech.edu/PAP/nonplsep.pdf}.}

\begin{lemma}\label{lem:path-small-components}
Let $J$ have $m$ vertices, and suppose its underlying simple graph has no $K_h$ minor.
For every integer $K\ge1$, there is a set $S$ such that
\[
 |S|\le\frac{4h^{3/2}m}{\sqrt K},
\]
and every component of $J-S$ has at most $K$ vertices.
\end{lemma}

\begin{proof}
Apply the separator theorem repeatedly to components with more than $K$ vertices.
Let $S$ contain all separators used in this recursion.
For a processed component of order $q$, charge $h^{3/2}/\sqrt q$ to each of its vertices.
The total charge covers its separator size.
For any fixed vertex, the orders of successive components containing it decrease by a factor of at least two.
Read that sequence in reverse order.
Its smallest order exceeds $K$, so the vertex's total charge is at most
\[
 \frac{h^{3/2}}{\sqrt K}\sum_{j\ge0}2^{-j/2}
 =\frac{(2+\sqrt2)h^{3/2}}{\sqrt K}
 <\frac{4h^{3/2}}{\sqrt K}.
\]
Sum over the vertices to obtain the stated bound.
\end{proof}

\begin{proof}[Proof of Theorem~\ref{thm:path-excluded-minor}]
Choose
\[
 \delta=\frac{\eta}{28+8\eta},\qquad
 K=\left\lceil\frac{16h^3}{\delta^2}\right\rceil,
 \qquad D=3K.
\]
Then $0<\delta<1/8$.
Apply Lemma~\ref{lem:path-route-reduction}, and put $J=\Gamma-V(A)$.
Lemma~\ref{lem:path-small-components} gives a set $S$ with
\[
 |S|\le\delta|V(J)|\le\delta(2b+\ell+1)
\]
such that every component of $J-S$ has at most $K$ vertices.
Deleting one vertex of degree at most five decreases rank by at most four.
Thus
\begin{equation}\label{eq:path-minor-surviving-rank}
 \beta(J-S)\ge (1-8\delta)b-(3+4\delta)\ell-4\delta.
\end{equation}

Every component of $J-S$ has maximum degree at most five and rank at most $2K$.
Choose one cycle from every cyclic component, obtaining $k$ vertex-disjoint cycles.
Their vertex sets lie outside $A$, and
\[
 2Kk\ge\beta(J-S).
\]
Each selected vertex has degree at most five in $\Gamma$ and uses two incidences in its cycle.
Every chosen cyclic region therefore has boundary at most $3K=D$.

Add a new two-edge terminal path to $\Gamma$, and mark its union with $A$.
The path injection and \eqref{eq:path-small-cuts} give
\[
 k\le\frac{(D-1)2^{D-2}}{\varepsilon}.
\]
For $k=0$, the same inequality holds trivially.
Combining these estimates with \eqref{eq:path-minor-surviving-rank} gives
\[
 (1-8\delta)b
 \le(3+4\delta)\ell+4\delta
 +\frac{2K(D-1)2^{D-2}}{\varepsilon}.
\]
The definition of $\delta$ gives
$(3+4\delta)/(1-8\delta)=3+\eta$.
Since $\varepsilon\le1/2$, one may take
\[
 C_{h,\eta}
 =\frac{2\delta+2K(D-1)2^{D-2}}{1-8\delta}.
\]
This proves \eqref{eq:path-minor-linear}.
Since $\ell\ge1$, it also gives
\[
 \ell\ge\frac{\varepsilon b}{(3+\eta)\varepsilon+C_{h,\eta}}.
\]
Finally, a proper minor-closed class omits some fixed finite graph $H$.
It then omits $K_{|V(H)|}$, because that complete graph contains $H$ as a minor.
\end{proof}

The constants above are deliberately explicit and are not optimized.
The proof establishes a structural consequence of the existing small-cut estimate and the classical separator theorem.
It does not provide a separator argument for arbitrary graphs.
